\section{Fuente de datos}

El archivo operativo es \code{data/raw/Spotify_Tracks_Dataset.csv}, proporcionado por la asignatura
y preservado sin modificaciones. Su estructura coincide con el conjunto publicado por Maharshi
Pandya en Kaggle \citep{pandya_dataset}. Esta coincidencia sustenta la descripción, pero la copia
local y su trazabilidad son la fuente efectiva del análisis.

El CSV original tiene 114.000 filas y 21 columnas. Incluye identificadores, artista, álbum, nombre
de pista, género, contenido explícito, duración y atributos acústicos. La variable objetivo
\code{popularity} está en escala 0--100. La documentación de Spotify la relaciona principalmente
con reproducciones y recencia, y el campo aparece actualmente como obsoleto en la API
\citep{spotify_track_api}. Se interpreta, por ello, como fotografía histórica y no como medición en
tiempo real.

La integridad de la copia raw se controla mediante SHA-256:

\begin{verbatim}
b202fa49909b2d5cef71a04b1d21243cfeb36414535f2ca9272aa646721177bd
\end{verbatim}

El diccionario completo está en \code{docs/data_dictionary.md}. En síntesis, los roles son:

\begin{table}[H]
\centering
\caption{Grupos de variables de la fuente.}
\label{tab:grupos-variables}
\begin{tabularx}{\textwidth}{p{3cm}X p{3.2cm}}
\toprule
Grupo & Variables representativas & Rol \\
\midrule
Identificación & \code{track_id} & Trazabilidad; no predictor \\
Contexto & \code{artists}, \code{album_name}, \code{track_name}, género & Segmentación y posible
alta cardinalidad \\
Audio & \code{danceability}, \code{energy}, \code{loudness}, \code{acousticness}, entre otras &
Predictores candidatos \\
Estructura musical & \code{key}, \code{mode}, \code{tempo}, \code{time_signature} & Predictores
categóricos o continuos según significado \\
Objetivo & \code{popularity} & Variable de respuesta histórica \\
Técnica & \code{Unnamed: 0} & Índice exportado, retirado \\
\bottomrule
\end{tabularx}
\end{table}

No se documenta la fecha exacta de extracción, la cobertura regional ni las condiciones completas
de muestreo. Esas ausencias limitan actualización, generalización temporal y representatividad.
